% Use the existing owner-derived source mechanism; do not edit frozen originals.
\OLSADeclareDocumentTextCorrection{OLP-0022}{FA0022-ONTO-INTRO}
{برای آغاز، ممکن است بخواهیم توابعی را بررسی کنیم که این ویژگی را دارند
که هر عضو هم‌دامنه مقداری از تابع است. چنین}
{برای آغاز، ممکن است بخواهیم توابعی را بررسی کنیم که هر عضو
هم‌دامنهٔ آن‌ها مقدار تابع برای دست‌کم یک ورودی باشد. چنین}
\OLSADeclareDocumentTextCorrection{OLP-0022}{FA0022-ONTO-CAPTION}
{\caption{تابع !!^a{surjective} هر !!{element} از
    هم‌دامنه را به‌عنوان یک مقدار دارد.}}
{\caption{در تابع !!^a{surjective}، هر !!{element} هم‌دامنه
    مقدار تابع برای دست‌کم یک ورودی است.}}
\OLSADeclareDocumentTextCorrection{OLP-0022}{FA0022-ONTO-DEFINITION}
{تابع $f \colon A \rightarrow B$ اگر و تنها اگر $B$
بردِ~$f$ نیز باشد، \emph{!!{surjective}} است؛ یعنی برای هر $y \in B$ دست‌کم}
{تابع $f \colon A \rightarrow B$ \emph{!!{surjective}} است اگر و تنها اگر $B$
همان بردِ~$f$ باشد؛ یعنی برای هر $y \in B$ دست‌کم}
\OLSADeclareDocumentTextCorrection{OLP-0022}{FA0022-INDUCED-ONTO}
{توجه کنید که هر تابعی !!a{surjection} \emph{القا می‌کند}. چراکه
با داشتن تابع $f \colon A \to B$، فرض کنید $f' \colon A \to \ran{f}$ به‌وسیلهٔ
$f'(x) = f(x)$ تعریف شود. چون $\ran{f}$ \emph{بنا بر تعریف} برابر است با
$\Setabs{f(x) \in B}{x \in A}$، تضمین می‌شود که این تابع $f'$
!!a{surjection} باشد}
{توجه کنید که از هر تابع می‌توان یک !!a{surjection} \emph{به دست آورد}.
برای تابع $f \colon A \to B$، تابع $f' \colon A \to \ran{f}$ را با
ضابطهٔ $f'(x) = f(x)$ تعریف می‌کنیم. چون $\ran{f}$ \emph{بنا بر تعریف} برابر است با
$\Setabs{f(x) \in B}{x \in A}$، این تابع $f'$
!!a{surjection} است.}
\OLSADeclareDocumentTextCorrection{OLP-0022}{FA0022-INJECTIVE-CAPTION}
{آرگومان متفاوت را به مقدار یکسان نگاشت نمی‌کند.}
{ورودی متفاوت را به مقدار یکسان نگاشت نمی‌کند.}
\OLSADeclareDocumentTextCorrection{OLP-0022}{FA0022-MEMBER-EZAFE}
{!!{element}sی از دامنهٔ $f$ هستند}
{!!{element}sی دامنهٔ $f$ هستند}
\OLSADeclareDocumentTextCorrection{OLP-0022}{FA0022-SUCCESSOR-CONSISTENCY}
{تابع پسین}
{تابع جانشین}
\OLSADeclareDocumentTextCorrection{OLP-0023}{FA0023-INTRODUCTION}
{تابعی که !!{element}s از~$A$ را به !!{element}s از~$B$
نگاشت می‌کند، آشکارا رابطه‌ای میان $A$ و~$B$ تعریف می‌کند؛ یعنی رابطه‌ای
که میان $x$ و $y$ برقرار است اگر و تنها اگر $f(x) = y$. در واقع، حتی ممکن است---
برای نمونه، اگر به تقلیل اجزای سازندهٔ ریاضیات علاقه‌مند باشیم---دست به
\emph{یکی‌انگاری} بزنیم و تابع~$f$ را با این رابطه، یعنی با مجموعه‌ای از
زوج‌ها، یکی بدانیم. آنگاه این پرسش پیش می‌آید:
کدام روابط به این شیوه توابع را تعریف می‌کنند؟}
{تابعی که !!{element}sی~$A$ را به !!{element}sی~$B$
نگاشت می‌کند، آشکارا رابطه‌ای میان $A$ و~$B$ تعریف می‌کند:
این رابطه میان $x$ و $y$ برقرار است اگر و تنها اگر $f(x) = y$.
حتی می‌توانیم---برای نمونه، اگر بخواهیم ریاضیات را بر مفاهیم بنیادیِ
کمتری بنا کنیم---تابع~$f$ را با این رابطه، یعنی با مجموعه‌ای از
زوج‌ها، \emph{یکی در نظر بگیریم}. آنگاه این پرسش پیش می‌آید:
کدام رابطه‌ها به این شیوه تابع تعریف می‌کنند؟}
\OLSADeclareDocumentTextCorrection{OLP-0023}{FA0023-ATTESTED-GRAPH}
{\begin{defn}[گراف یک تابع] فرض کنید $f\colon A \to B$ یک تابع باشد.
\emph{گرافِ}~$f$ رابطهٔ $R_f \subseteq A \times B$ است که}
{\begin{defn}[نمودار یک تابع] فرض کنید $f\colon A \to B$ یک تابع باشد.
\emph{نمودارِ}~$f$ رابطهٔ $R_f \subseteq A \times B$ است که}
\OLSADeclareDocumentTextCorrection{OLP-0023}{FA0023-EXTENSIONALITY}
{گراف یک تابع بنا بر اصل گسترش به‌طور یکتا تعیین می‌شود.
افزون بر این، اصل گسترش (برای مجموعه‌ها) بی‌درنگ
اصل ضمنی گسترش برای توابع را موجه می‌سازد؛
بنابر این اصل، اگر $f$ و~$g$ دامنه و هم‌دامنهٔ یکسانی داشته باشند، آنگاه اگر
برای همهٔ آرگومان‌ها مقادیر یکسانی داشته باشند، یکسان‌اند.}
{نمودار یک تابع بنا بر اصل گسترش به‌طور یکتا تعیین می‌شود.
از اصل گسترش برای مجموعه‌ها، اصل گسترش برای توابع نیز نتیجه می‌شود؛
همان اصلی که تا اینجا به‌طور ضمنی به کار برده‌ایم:
اگر $f$ و~$g$ دامنه و هم‌دامنهٔ یکسانی داشته باشند و برای هر
ورودی مقدار یکسانی بدهند، دو تابع برابرند.}
\OLSADeclareDocumentTextCorrection{OLP-0023}{FA0023-FUNCTIONAL-DOMAIN}
{به همین ترتیب، اگر رابطه‌ای «تابع‌وار» باشد، آنگاه گراف یک تابع است.}
{به همین ترتیب، رابطه‌ای که هر مؤلفهٔ نخست را حداکثر به یک مؤلفهٔ دوم مربوط کند («تابع‌وار» باشد)، نمودار تابعی بر دامنهٔ همان رابطه است.}
\OLSADeclareDocumentTextCorrection{OLP-0023}{FA0023-PROPOSITION-GRAPH}
{آنگاه $R$ گراف تابع $f\colon A \to B$ است که به‌وسیلهٔ}
{آنگاه $R$ نمودار تابع $f\colon A \to B$ است که به‌وسیلهٔ}
\OLSADeclareDocumentTextCorrection{OLP-0023}{FA0023-EXISTENCE-EXPLICIT}
{فرض کنید یک $y$ وجود دارد چنان‌که $Rxy$.}
{به‌ازای هر ورودی در دامنه، شرط دوم وجود یک $y$ را تضمین می‌کند چنان‌که $Rxy$.}
\OLSADeclareDocumentTextCorrection{OLP-0023}{FA0023-REPRESENTATION}
{هر تابع $f\colon A \to B$ گرافی دارد؛ یعنی رابطه‌ای روی $A
\times B$ که با $f(x) = y$ تعریف می‌شود. از سوی دیگر، هر رابطهٔ~$R
\subseteq A \times B$ که ویژگی‌های داده‌شده در
\olref{prop:graph-function} را داشته باشد، گراف تابعی~$f \colon A \to
B$ است. به‌سبب این پیوند نزدیک میان توابع و
گراف‌هایشان، می‌توانیم یک تابع را صرفاً گراف آن در نظر بگیریم. به
بیان دیگر، توابع را می‌توان با روابطی معین، یعنی با
مجموعه‌هایی معین از چندتایی‌ها، یکی انگاشت. \oliflabeldef{sfr:rel:ref:sec}{بااین‌حال،
توجه کنید که روح این «یکی‌انگاری» همان است که در
\olref[sfr][rel][ref]{sec} آمده است: این ادعایی دربارهٔ متافیزیک
توابع نیست، بلکه مشاهده‌ای است مبنی بر اینکه مناسب است توابع را همچون
مجموعه‌هایی معین \emph{در نظر بگیریم}. یکی از دلایل این مناسب‌بودن آن است
که اکنون می‌توانیم}{اکنون می‌توانیم} عملیات مشابهی را روی
توابع انجام دهیم که روی روابط انجام دادیم (نگاه کنید به
\olref[sfr][rel][ops]{sec}). به‌ویژه:}
{هر تابع $f\colon A \to B$ نموداری دارد؛ یعنی رابطه‌ای به‌صورت زیرمجموعه‌ای از $A
\times B$ که با $f(x) = y$ تعریف می‌شود. از سوی دیگر، هر رابطهٔ~$R
\subseteq A \times B$ که ویژگی‌های داده‌شده در
\olref{prop:graph-function} را داشته باشد، نمودار تابعی~$f \colon A \to
B$ است. به‌سبب این پیوند نزدیک میان توابع و
نمودارهایشان، می‌توانیم تابع را همان نمودار آن در نظر بگیریم. به
بیان دیگر، می‌توان توابع را با روابطی معین، یعنی با
مجموعه‌هایی معین از چندتایی‌ها، یکی در نظر گرفت. \oliflabeldef{sfr:rel:ref:sec}{بااین‌حال،
منظور از این «یکی در نظر گرفتن» همان است که در
\olref[sfr][rel][ref]{sec} توضیح دادیم: ادعا نمی‌کنیم که ماهیت توابع چیست؛
فقط می‌گوییم که برای کار ریاضی مفید است توابع را مجموعه‌هایی معین
\emph{در نظر بگیریم}. یکی از فایده‌های این روش آن است
که اکنون می‌توانیم}{اکنون می‌توانیم} همان‌گونه که روی روابط عملیات انجام دادیم،
روی توابع نیز عملیات مشابهی انجام دهیم (نگاه کنید به
\olref[sfr][rel][ops]{sec}). به‌ویژه:}
\OLSADeclareDocumentTextCorrection{OLP-0023}{FA0023-RESTRICTION-CONVENTION}
{\oliflabeldef{sfr:rel:ops:sec}{این مفاهیم با توجه به تعریف‌های
\olref[sfr][rel][ops]{sec} و یکی‌انگاری توابع با روابط، دقیقاً
همان‌اند که انتظار می‌رود. اما دو عمل دیگر---وارون‌ها و حاصل‌ضرب‌های نسبی---به
جزئیات اندکی بیشتر نیاز دارند. آن جزئیات را در
\olref[inv]{sec} و
\olref[cmp]{sec} ارائه خواهیم کرد.}{}}
{\oliflabeldef{sfr:rel:ops:sec}{این مفاهیم را می‌توان با تعریف‌های
\olref[sfr][rel][ops]{sec} و در نظر گرفتن توابع به‌صورت روابط مقایسه کرد.
\emph{یادداشت ویراستاری:} تعریف تصویر با تعریف قبلی سازگار است؛
اما قرارداد تحدید در اینجا تفاوت دارد. در تحدید تابع، فقط ورودی‌ها
به زیرمجموعهٔ دامنه محدود می‌شوند و هم‌دامنه تغییر نمی‌کند؛
در تعریف پیشینِ تحدید رابطه، هر دو مؤلفه به مجموعهٔ تعیین‌شده محدود می‌شدند.
دو عمل دیگر---وارون‌ها و حاصل‌ضرب‌های نسبی---به
جزئیات اندکی بیشتر نیاز دارند. آن جزئیات را در
\olref[inv]{sec} و
\olref[cmp]{sec} ارائه خواهیم کرد.}{}}
\OLSADeclareSetup{OLP-0022}{\olsa_source_correction_install_document_text:}
\OLSADeclareSetup{OLP-0023}{\olsa_source_correction_install_document_text:}
